% List-blind (off-market) part of the team report. Paste into report.tex, or % List-blind (off-market) part of the team report. Paste into report.tex, or % List-blind (off-market) part of the team report. Paste into report.tex, or % List-blind (off-market) part of the team report. Paste into report.tex, or \input{list_blind_section}.
% Every number comes from our own validation: `python run.py backtest`, `validate` and `report`.

\section{Off-market (list-blind) model}

\textbf{Task and data.} We value a home with no list price. We train on the 150,000 sales from January 2021
to May 2026 and predict the 3,974 off-market test homes, which sold June to September 2026. We use every
column the test file has. The columns that only exist after a sale (sale price, sale date, days on market,
the external AVM) are dropped from every home we predict, and the list price is never used. Days on
market from \emph{past} sales is used only to guess when a test home is likely to sell.

\textbf{Features (52 per home).}
(i) \emph{Market index:} a ridge regression of log price on community $\times$ type, size, bedrooms,
washrooms, parking and one dummy per sale month. The month coefficients give a monthly market level, net of
what sold. It peaks in March 2022, is about 25\% lower by December 2025, and dips from spring into summer
every year. Every training price is divided by its month's level, so the trees learn what a home is worth
at an average market level and the index alone carries time.
(ii) \emph{The home:} size (the midpoint when given as a range, plus a flag), bedrooms, washrooms,
parking, garage, lot frontage, depth and area, type, style, municipality, community, latitude and longitude.
(iii) \emph{Same building:} earlier sales at the exact same coordinates, as a count and an average
market-adjusted price per sq ft.
(iv) \emph{Nearest sales:} the median adjusted price per sq ft of the 10 nearest earlier sales of the
same type, times this home's size, and their average distance.
(v) \emph{Listing text:} a ridge regression on TF-IDF word and word-pair counts predicting how far a
home's price per sq ft sits above its community-and-type average, plus 30 SVD topics. The strongest words
are size cues, then condition (``needs'' and ``as is'' down, ``renovated'' up) and some street names.
\textbf{No leakage:} a training sale only sees sales from earlier months, building statistics exclude
the sale itself, and training rows get out-of-fold text scores.

\textbf{Model.} The target is the market-adjusted log price, and a sale's weight halves every 12 months.
The value is the average of four histogram gradient-boosting models (scikit-learn): two with 2,000 trees at
learning rate 0.025 and two with 800 trees at 0.05. Every leaf holds at least 20 sales, 30\% of features
are tried at each split, and each model has its own seed. On the 2024 summer one tree alone misses by 11.1\%,
the average of four by 6.85\%.
The \emph{market level} for a test home averages the projected index over its likely sale dates: its
listing date plus 100 typical times on market from past sales, kept only if inside June to September.
Months past May are projected as the last level plus the median change for that calendar month in earlier
years, so one unusual year cannot drag the projection.
The \emph{local fix}: five-fold out-of-fold predictions give an honest miss for every training sale. For
a new home we average the misses of its 5, 10, 20 and 40 nearest earlier sales of the same type, and
combine the four averages with weights fitted per property type by ridge regression.
The estimate is $\exp(\textrm{value} + \textrm{market level} + \textrm{local fix})$, so the three log
parts multiply in dollars. For one detached house in the 2024 backtest: \$898,771 $\times$ 0.990 (June
market) $\times$ 1.121 (its area sells above the model's guesses) $=$ \$997,309. It sold for \$994,600.

\textbf{How it differs from the list-aware model, and why.} The two tracks answer different questions,
so they are built differently. The list price already carries most of what the market knows: the current
price level, and the home's condition and finish as the seller's agent sees them. The list-aware model
therefore only predicts the gap between sale and list price (CatBoost on log(sale / list), blended with a
ridge regression, separately for condos and freeholds), and needs no market index or comparable sales.
The off-market model may not use the list price, so it has to rebuild the value from scratch: a market
index to bring every past sale to the same price level and project it into the summer, nearest-sale and
same-building features to stand in for what an agent would know, the listing text as a hint about
condition, and the local fix for what is still missing. That is why its error is about twice as large.

\textbf{Validation.} \emph{Walk-forward} (the figure we report): each month from December 2025 to May 2026
is predicted by a model trained only on sales before it, six full refits. \emph{Summer backtests} copy the
real test: train on sales before June 1, predict June to September. Every modelling choice was made on the
2023 and 2024 summers. 2025 was held back; it was checked once at the end and again after the last two
changes, so we treat it as a confirmation rather than a pristine holdout. We also ran 2021 and 2022 with the
final model. We never used a random split.

\begin{table}[h]
\centering\small
\begin{tabular}{lccccc}
\toprule
Summer (train before June 1) & 2021 & 2022 & 2023 & 2024 & 2025 (held back) \\
\midrule
MdAPE / MAPE \% & 7.56 / 10.07 & 8.30 / 11.60 & 6.62 / 8.88 & 6.64 / 8.96 & 6.84 / 9.08 \\
Median signed error \% & +0.4 & +5.1 & +0.9 & +1.5 & +0.9 \\
\bottomrule
\end{tabular}
\caption{Off-market summer backtests. 2021 had only five months of training data; 2022 is the year
prices fell fast after the spring peak.}
\end{table}

\textbf{Results.} On walk-forward validation the off-market MdAPE is \textbf{7.31\%} (bootstrap 95\%
interval 7.14--7.50; MAPE 9.44\%), against 17.7\% for the organisers' simple baseline and 20.8\% for the
starter (median of the 10 nearest sales). Validation is higher than the summers because March and April 2026
were hard months (7.9\% and 7.7\%) as prices rose quickly that spring; the other four months were 6.8 to
7.2\%. The hidden test is a summer, so we expect about \textbf{6.7 to 7.0\%} unless the summer market turns
sharply, as in 2022. Across all 54,991 out-of-sample predictions, 64\% land within 10\% of the sale price
and 89\% within 20\%.

\begin{figure}[h]
\centering
\includegraphics[width=0.44\linewidth]{slides/failures_type.png}\hspace{0.04\linewidth}
\includegraphics[width=0.44\linewidth]{slides/failures_decile.png}
\caption{Off-market typical miss on the 2025 summer by property type (homes in brackets) and by tenth of
predicted price.}
\end{figure}

\textbf{Where it fails, and why.} \emph{Type:} condos miss by about 5.5\%, houses by 8.4 to 8.7\%.
Condos in one building are near-copies, so building history pins them down; two houses with the same size,
rooms and street can differ in age, renovation, lot and finish, none of which the data records.
\emph{Price:} the cheapest tenth misses by 5.1\%, the dearest by 10.5\%, since expensive homes are rarer,
more varied and mostly houses. \emph{Neighbourhood:} across the 97 communities with at least 40 sales in
2025, the typical miss ranges from 3.9\% (Cedarglen, condos) to 13.5\% (Lindencliffe, 73\% houses);
the middle half is 6.1 to 8.0\%. It is worse where more sales are houses (correlation 0.39) and where prices
are higher (0.43). \emph{Time:} estimates lean about 2\% high in July and August, when the summer dip runs
deeper than projected, and 5\% high in 2022, when prices fell after the spring peak.
\emph{Honest uncertainty:} an 80\% range for each estimate, from the 10th to 90th percentile of the
2023--24 misses for its property type, covered 79.5\% of the 2025 homes, with a typical width of 22.6\% of
the price (about $\pm$11\%). The bootstrap intervals treat homes as independent; homes sold in the same
month share the same market surprise, so the real uncertainty is somewhat wider.

\textbf{What did not help} (all judged on the 2023 and 2024 summers): separate indexes for houses and
condos (the projection noise outweighed the bias it fixed); following community price trends (changes in a
community's premium have autocorrelation $-0.41$, so they reverse rather than persist); blending in a random
forest (9.2\% on its own); separate tree models for houses and condos (within noise); a fixed summer bias correction (the best shift varied from 0 to 1.75\% by year).
\textbf{Limits:} do not trust it for expensive or unusual houses, for thin groups such as co-ops (50 test
homes in 2025), or in a summer when the market turns sharply. Data on age, renovation, floor level and condo
fees would help most.
.
% Every number comes from our own validation: `python run.py backtest`, `validate` and `report`.

\section{Off-market (list-blind) model}

\textbf{Task and data.} We value a home with no list price. We train on the 150,000 sales from January 2021
to May 2026 and predict the 3,974 off-market test homes, which sold June to September 2026. We use every
column the test file has. The columns that only exist after a sale (sale price, sale date, days on market,
the external AVM) are dropped from every home we predict, and the list price is never used. Days on
market from \emph{past} sales is used only to guess when a test home is likely to sell.

\textbf{Features (52 per home).}
(i) \emph{Market index:} a ridge regression of log price on community $\times$ type, size, bedrooms,
washrooms, parking and one dummy per sale month. The month coefficients give a monthly market level, net of
what sold. It peaks in March 2022, is about 25\% lower by December 2025, and dips from spring into summer
every year. Every training price is divided by its month's level, so the trees learn what a home is worth
at an average market level and the index alone carries time.
(ii) \emph{The home:} size (the midpoint when given as a range, plus a flag), bedrooms, washrooms,
parking, garage, lot frontage, depth and area, type, style, municipality, community, latitude and longitude.
(iii) \emph{Same building:} earlier sales at the exact same coordinates, as a count and an average
market-adjusted price per sq ft.
(iv) \emph{Nearest sales:} the median adjusted price per sq ft of the 10 nearest earlier sales of the
same type, times this home's size, and their average distance.
(v) \emph{Listing text:} a ridge regression on TF-IDF word and word-pair counts predicting how far a
home's price per sq ft sits above its community-and-type average, plus 30 SVD topics. The strongest words
are size cues, then condition (``needs'' and ``as is'' down, ``renovated'' up) and some street names.
\textbf{No leakage:} a training sale only sees sales from earlier months, building statistics exclude
the sale itself, and training rows get out-of-fold text scores.

\textbf{Model.} The target is the market-adjusted log price, and a sale's weight halves every 12 months.
The value is the average of four histogram gradient-boosting models (scikit-learn): two with 2,000 trees at
learning rate 0.025 and two with 800 trees at 0.05. Every leaf holds at least 20 sales, 30\% of features
are tried at each split, and each model has its own seed. On the 2024 summer one tree alone misses by 11.1\%,
the average of four by 6.85\%.
The \emph{market level} for a test home averages the projected index over its likely sale dates: its
listing date plus 100 typical times on market from past sales, kept only if inside June to September.
Months past May are projected as the last level plus the median change for that calendar month in earlier
years, so one unusual year cannot drag the projection.
The \emph{local fix}: five-fold out-of-fold predictions give an honest miss for every training sale. For
a new home we average the misses of its 5, 10, 20 and 40 nearest earlier sales of the same type, and
combine the four averages with weights fitted per property type by ridge regression.
The estimate is $\exp(\textrm{value} + \textrm{market level} + \textrm{local fix})$, so the three log
parts multiply in dollars. For one detached house in the 2024 backtest: \$898,771 $\times$ 0.990 (June
market) $\times$ 1.121 (its area sells above the model's guesses) $=$ \$997,309. It sold for \$994,600.

\textbf{How it differs from the list-aware model, and why.} The two tracks answer different questions,
so they are built differently. The list price already carries most of what the market knows: the current
price level, and the home's condition and finish as the seller's agent sees them. The list-aware model
therefore only predicts the gap between sale and list price (CatBoost on log(sale / list), blended with a
ridge regression, separately for condos and freeholds), and needs no market index or comparable sales.
The off-market model may not use the list price, so it has to rebuild the value from scratch: a market
index to bring every past sale to the same price level and project it into the summer, nearest-sale and
same-building features to stand in for what an agent would know, the listing text as a hint about
condition, and the local fix for what is still missing. That is why its error is about twice as large.

\textbf{Validation.} \emph{Walk-forward} (the figure we report): each month from December 2025 to May 2026
is predicted by a model trained only on sales before it, six full refits. \emph{Summer backtests} copy the
real test: train on sales before June 1, predict June to September. Every modelling choice was made on the
2023 and 2024 summers. 2025 was held back; it was checked once at the end and again after the last two
changes, so we treat it as a confirmation rather than a pristine holdout. We also ran 2021 and 2022 with the
final model. We never used a random split.

\begin{table}[h]
\centering\small
\begin{tabular}{lccccc}
\toprule
Summer (train before June 1) & 2021 & 2022 & 2023 & 2024 & 2025 (held back) \\
\midrule
MdAPE / MAPE \% & 7.56 / 10.07 & 8.30 / 11.60 & 6.62 / 8.88 & 6.64 / 8.96 & 6.84 / 9.08 \\
Median signed error \% & +0.4 & +5.1 & +0.9 & +1.5 & +0.9 \\
\bottomrule
\end{tabular}
\caption{Off-market summer backtests. 2021 had only five months of training data; 2022 is the year
prices fell fast after the spring peak.}
\end{table}

\textbf{Results.} On walk-forward validation the off-market MdAPE is \textbf{7.31\%} (bootstrap 95\%
interval 7.14--7.50; MAPE 9.44\%), against 17.7\% for the organisers' simple baseline and 20.8\% for the
starter (median of the 10 nearest sales). Validation is higher than the summers because March and April 2026
were hard months (7.9\% and 7.7\%) as prices rose quickly that spring; the other four months were 6.8 to
7.2\%. The hidden test is a summer, so we expect about \textbf{6.7 to 7.0\%} unless the summer market turns
sharply, as in 2022. Across all 54,991 out-of-sample predictions, 64\% land within 10\% of the sale price
and 89\% within 20\%.

\begin{figure}[h]
\centering
\includegraphics[width=0.44\linewidth]{slides/failures_type.png}\hspace{0.04\linewidth}
\includegraphics[width=0.44\linewidth]{slides/failures_decile.png}
\caption{Off-market typical miss on the 2025 summer by property type (homes in brackets) and by tenth of
predicted price.}
\end{figure}

\textbf{Where it fails, and why.} \emph{Type:} condos miss by about 5.5\%, houses by 8.4 to 8.7\%.
Condos in one building are near-copies, so building history pins them down; two houses with the same size,
rooms and street can differ in age, renovation, lot and finish, none of which the data records.
\emph{Price:} the cheapest tenth misses by 5.1\%, the dearest by 10.5\%, since expensive homes are rarer,
more varied and mostly houses. \emph{Neighbourhood:} across the 97 communities with at least 40 sales in
2025, the typical miss ranges from 3.9\% (Cedarglen, condos) to 13.5\% (Lindencliffe, 73\% houses);
the middle half is 6.1 to 8.0\%. It is worse where more sales are houses (correlation 0.39) and where prices
are higher (0.43). \emph{Time:} estimates lean about 2\% high in July and August, when the summer dip runs
deeper than projected, and 5\% high in 2022, when prices fell after the spring peak.
\emph{Honest uncertainty:} an 80\% range for each estimate, from the 10th to 90th percentile of the
2023--24 misses for its property type, covered 79.5\% of the 2025 homes, with a typical width of 22.6\% of
the price (about $\pm$11\%). The bootstrap intervals treat homes as independent; homes sold in the same
month share the same market surprise, so the real uncertainty is somewhat wider.

\textbf{What did not help} (all judged on the 2023 and 2024 summers): separate indexes for houses and
condos (the projection noise outweighed the bias it fixed); following community price trends (changes in a
community's premium have autocorrelation $-0.41$, so they reverse rather than persist); blending in a random
forest (9.2\% on its own); separate tree models for houses and condos (within noise); a fixed summer bias correction (the best shift varied from 0 to 1.75\% by year).
\textbf{Limits:} do not trust it for expensive or unusual houses, for thin groups such as co-ops (50 test
homes in 2025), or in a summer when the market turns sharply. Data on age, renovation, floor level and condo
fees would help most.
.
% Every number comes from our own validation: `python run.py backtest`, `validate` and `report`.

\section{Off-market (list-blind) model}

\textbf{Task and data.} We value a home with no list price. We train on the 150,000 sales from January 2021
to May 2026 and predict the 3,974 off-market test homes, which sold June to September 2026. We use every
column the test file has. The columns that only exist after a sale (sale price, sale date, days on market,
the external AVM) are dropped from every home we predict, and the list price is never used. Days on
market from \emph{past} sales is used only to guess when a test home is likely to sell.

\textbf{Features (52 per home).}
(i) \emph{Market index:} a ridge regression of log price on community $\times$ type, size, bedrooms,
washrooms, parking and one dummy per sale month. The month coefficients give a monthly market level, net of
what sold. It peaks in March 2022, is about 25\% lower by December 2025, and dips from spring into summer
every year. Every training price is divided by its month's level, so the trees learn what a home is worth
at an average market level and the index alone carries time.
(ii) \emph{The home:} size (the midpoint when given as a range, plus a flag), bedrooms, washrooms,
parking, garage, lot frontage, depth and area, type, style, municipality, community, latitude and longitude.
(iii) \emph{Same building:} earlier sales at the exact same coordinates, as a count and an average
market-adjusted price per sq ft.
(iv) \emph{Nearest sales:} the median adjusted price per sq ft of the 10 nearest earlier sales of the
same type, times this home's size, and their average distance.
(v) \emph{Listing text:} a ridge regression on TF-IDF word and word-pair counts predicting how far a
home's price per sq ft sits above its community-and-type average, plus 30 SVD topics. The strongest words
are size cues, then condition (``needs'' and ``as is'' down, ``renovated'' up) and some street names.
\textbf{No leakage:} a training sale only sees sales from earlier months, building statistics exclude
the sale itself, and training rows get out-of-fold text scores.

\textbf{Model.} The target is the market-adjusted log price, and a sale's weight halves every 12 months.
The value is the average of four histogram gradient-boosting models (scikit-learn): two with 2,000 trees at
learning rate 0.025 and two with 800 trees at 0.05. Every leaf holds at least 20 sales, 30\% of features
are tried at each split, and each model has its own seed. On the 2024 summer one tree alone misses by 11.1\%,
the average of four by 6.85\%.
The \emph{market level} for a test home averages the projected index over its likely sale dates: its
listing date plus 100 typical times on market from past sales, kept only if inside June to September.
Months past May are projected as the last level plus the median change for that calendar month in earlier
years, so one unusual year cannot drag the projection.
The \emph{local fix}: five-fold out-of-fold predictions give an honest miss for every training sale. For
a new home we average the misses of its 5, 10, 20 and 40 nearest earlier sales of the same type, and
combine the four averages with weights fitted per property type by ridge regression.
The estimate is $\exp(\textrm{value} + \textrm{market level} + \textrm{local fix})$, so the three log
parts multiply in dollars. For one detached house in the 2024 backtest: \$898,771 $\times$ 0.990 (June
market) $\times$ 1.121 (its area sells above the model's guesses) $=$ \$997,309. It sold for \$994,600.

\textbf{How it differs from the list-aware model, and why.} The two tracks answer different questions,
so they are built differently. The list price already carries most of what the market knows: the current
price level, and the home's condition and finish as the seller's agent sees them. The list-aware model
therefore only predicts the gap between sale and list price (CatBoost on log(sale / list), blended with a
ridge regression, separately for condos and freeholds), and needs no market index or comparable sales.
The off-market model may not use the list price, so it has to rebuild the value from scratch: a market
index to bring every past sale to the same price level and project it into the summer, nearest-sale and
same-building features to stand in for what an agent would know, the listing text as a hint about
condition, and the local fix for what is still missing. That is why its error is about twice as large.

\textbf{Validation.} \emph{Walk-forward} (the figure we report): each month from December 2025 to May 2026
is predicted by a model trained only on sales before it, six full refits. \emph{Summer backtests} copy the
real test: train on sales before June 1, predict June to September. Every modelling choice was made on the
2023 and 2024 summers. 2025 was held back; it was checked once at the end and again after the last two
changes, so we treat it as a confirmation rather than a pristine holdout. We also ran 2021 and 2022 with the
final model. We never used a random split.

\begin{table}[h]
\centering\small
\begin{tabular}{lccccc}
\toprule
Summer (train before June 1) & 2021 & 2022 & 2023 & 2024 & 2025 (held back) \\
\midrule
MdAPE / MAPE \% & 7.56 / 10.07 & 8.30 / 11.60 & 6.62 / 8.88 & 6.64 / 8.96 & 6.84 / 9.08 \\
Median signed error \% & +0.4 & +5.1 & +0.9 & +1.5 & +0.9 \\
\bottomrule
\end{tabular}
\caption{Off-market summer backtests. 2021 had only five months of training data; 2022 is the year
prices fell fast after the spring peak.}
\end{table}

\textbf{Results.} On walk-forward validation the off-market MdAPE is \textbf{7.31\%} (bootstrap 95\%
interval 7.14--7.50; MAPE 9.44\%), against 17.7\% for the organisers' simple baseline and 20.8\% for the
starter (median of the 10 nearest sales). Validation is higher than the summers because March and April 2026
were hard months (7.9\% and 7.7\%) as prices rose quickly that spring; the other four months were 6.8 to
7.2\%. The hidden test is a summer, so we expect about \textbf{6.7 to 7.0\%} unless the summer market turns
sharply, as in 2022. Across all 54,991 out-of-sample predictions, 64\% land within 10\% of the sale price
and 89\% within 20\%.

\begin{figure}[h]
\centering
\includegraphics[width=0.44\linewidth]{slides/failures_type.png}\hspace{0.04\linewidth}
\includegraphics[width=0.44\linewidth]{slides/failures_decile.png}
\caption{Off-market typical miss on the 2025 summer by property type (homes in brackets) and by tenth of
predicted price.}
\end{figure}

\textbf{Where it fails, and why.} \emph{Type:} condos miss by about 5.5\%, houses by 8.4 to 8.7\%.
Condos in one building are near-copies, so building history pins them down; two houses with the same size,
rooms and street can differ in age, renovation, lot and finish, none of which the data records.
\emph{Price:} the cheapest tenth misses by 5.1\%, the dearest by 10.5\%, since expensive homes are rarer,
more varied and mostly houses. \emph{Neighbourhood:} across the 97 communities with at least 40 sales in
2025, the typical miss ranges from 3.9\% (Cedarglen, condos) to 13.5\% (Lindencliffe, 73\% houses);
the middle half is 6.1 to 8.0\%. It is worse where more sales are houses (correlation 0.39) and where prices
are higher (0.43). \emph{Time:} estimates lean about 2\% high in July and August, when the summer dip runs
deeper than projected, and 5\% high in 2022, when prices fell after the spring peak.
\emph{Honest uncertainty:} an 80\% range for each estimate, from the 10th to 90th percentile of the
2023--24 misses for its property type, covered 79.5\% of the 2025 homes, with a typical width of 22.6\% of
the price (about $\pm$11\%). The bootstrap intervals treat homes as independent; homes sold in the same
month share the same market surprise, so the real uncertainty is somewhat wider.

\textbf{What did not help} (all judged on the 2023 and 2024 summers): separate indexes for houses and
condos (the projection noise outweighed the bias it fixed); following community price trends (changes in a
community's premium have autocorrelation $-0.41$, so they reverse rather than persist); blending in a random
forest (9.2\% on its own); separate tree models for houses and condos (within noise); a fixed summer bias correction (the best shift varied from 0 to 1.75\% by year).
\textbf{Limits:} do not trust it for expensive or unusual houses, for thin groups such as co-ops (50 test
homes in 2025), or in a summer when the market turns sharply. Data on age, renovation, floor level and condo
fees would help most.
.
% Every number comes from our own validation: `python run.py backtest`, `validate` and `report`.

\section{Off-market (list-blind) model}

\textbf{Task and data.} We value a home with no list price. We train on the 150,000 sales from January 2021
to May 2026 and predict the 3,974 off-market test homes, which sold June to September 2026. We use every
column the test file has. The columns that only exist after a sale (sale price, sale date, days on market,
the external AVM) are dropped from every home we predict, and the list price is never used. Days on
market from \emph{past} sales is used only to guess when a test home is likely to sell.

\textbf{Features (52 per home).}
(i) \emph{Market index:} a ridge regression of log price on community $\times$ type, size, bedrooms,
washrooms, parking and one dummy per sale month. The month coefficients give a monthly market level, net of
what sold. It peaks in March 2022, is about 25\% lower by December 2025, and dips from spring into summer
every year. Every training price is divided by its month's level, so the trees learn what a home is worth
at an average market level and the index alone carries time.
(ii) \emph{The home:} size (the midpoint when given as a range, plus a flag), bedrooms, washrooms,
parking, garage, lot frontage, depth and area, type, style, municipality, community, latitude and longitude.
(iii) \emph{Same building:} earlier sales at the exact same coordinates, as a count and an average
market-adjusted price per sq ft.
(iv) \emph{Nearest sales:} the median adjusted price per sq ft of the 10 nearest earlier sales of the
same type, times this home's size, and their average distance.
(v) \emph{Listing text:} a ridge regression on TF-IDF word and word-pair counts predicting how far a
home's price per sq ft sits above its community-and-type average, plus 30 SVD topics. The strongest words
are size cues, then condition (``needs'' and ``as is'' down, ``renovated'' up) and some street names.
\textbf{No leakage:} a training sale only sees sales from earlier months, building statistics exclude
the sale itself, and training rows get out-of-fold text scores.

\textbf{Model.} The target is the market-adjusted log price, and a sale's weight halves every 12 months.
The value is the average of four histogram gradient-boosting models (scikit-learn): two with 2,000 trees at
learning rate 0.025 and two with 800 trees at 0.05. Every leaf holds at least 20 sales, 30\% of features
are tried at each split, and each model has its own seed. On the 2024 summer one tree alone misses by 11.1\%,
the average of four by 6.85\%.
The \emph{market level} for a test home averages the projected index over its likely sale dates: its
listing date plus 100 typical times on market from past sales, kept only if inside June to September.
Months past May are projected as the last level plus the median change for that calendar month in earlier
years, so one unusual year cannot drag the projection.
The \emph{local fix}: five-fold out-of-fold predictions give an honest miss for every training sale. For
a new home we average the misses of its 5, 10, 20 and 40 nearest earlier sales of the same type, and
combine the four averages with weights fitted per property type by ridge regression.
The estimate is $\exp(\textrm{value} + \textrm{market level} + \textrm{local fix})$, so the three log
parts multiply in dollars. For one detached house in the 2024 backtest: \$898,771 $\times$ 0.990 (June
market) $\times$ 1.121 (its area sells above the model's guesses) $=$ \$997,309. It sold for \$994,600.

\textbf{How it differs from the list-aware model, and why.} The two tracks answer different questions,
so they are built differently. The list price already carries most of what the market knows: the current
price level, and the home's condition and finish as the seller's agent sees them. The list-aware model
therefore only predicts the gap between sale and list price (CatBoost on log(sale / list), blended with a
ridge regression, separately for condos and freeholds), and needs no market index or comparable sales.
The off-market model may not use the list price, so it has to rebuild the value from scratch: a market
index to bring every past sale to the same price level and project it into the summer, nearest-sale and
same-building features to stand in for what an agent would know, the listing text as a hint about
condition, and the local fix for what is still missing. That is why its error is about twice as large.

\textbf{Validation.} \emph{Walk-forward} (the figure we report): each month from December 2025 to May 2026
is predicted by a model trained only on sales before it, six full refits. \emph{Summer backtests} copy the
real test: train on sales before June 1, predict June to September. Every modelling choice was made on the
2023 and 2024 summers. 2025 was held back; it was checked once at the end and again after the last two
changes, so we treat it as a confirmation rather than a pristine holdout. We also ran 2021 and 2022 with the
final model. We never used a random split.

\begin{table}[h]
\centering\small
\begin{tabular}{lccccc}
\toprule
Summer (train before June 1) & 2021 & 2022 & 2023 & 2024 & 2025 (held back) \\
\midrule
MdAPE / MAPE \% & 7.56 / 10.07 & 8.30 / 11.60 & 6.62 / 8.88 & 6.64 / 8.96 & 6.84 / 9.08 \\
Median signed error \% & +0.4 & +5.1 & +0.9 & +1.5 & +0.9 \\
\bottomrule
\end{tabular}
\caption{Off-market summer backtests. 2021 had only five months of training data; 2022 is the year
prices fell fast after the spring peak.}
\end{table}

\textbf{Results.} On walk-forward validation the off-market MdAPE is \textbf{7.31\%} (bootstrap 95\%
interval 7.14--7.50; MAPE 9.44\%), against 17.7\% for the organisers' simple baseline and 20.8\% for the
starter (median of the 10 nearest sales). Validation is higher than the summers because March and April 2026
were hard months (7.9\% and 7.7\%) as prices rose quickly that spring; the other four months were 6.8 to
7.2\%. The hidden test is a summer, so we expect about \textbf{6.7 to 7.0\%} unless the summer market turns
sharply, as in 2022. Across all 54,991 out-of-sample predictions, 64\% land within 10\% of the sale price
and 89\% within 20\%.

\begin{figure}[h]
\centering
\includegraphics[width=0.44\linewidth]{slides/failures_type.png}\hspace{0.04\linewidth}
\includegraphics[width=0.44\linewidth]{slides/failures_decile.png}
\caption{Off-market typical miss on the 2025 summer by property type (homes in brackets) and by tenth of
predicted price.}
\end{figure}

\textbf{Where it fails, and why.} \emph{Type:} condos miss by about 5.5\%, houses by 8.4 to 8.7\%.
Condos in one building are near-copies, so building history pins them down; two houses with the same size,
rooms and street can differ in age, renovation, lot and finish, none of which the data records.
\emph{Price:} the cheapest tenth misses by 5.1\%, the dearest by 10.5\%, since expensive homes are rarer,
more varied and mostly houses. \emph{Neighbourhood:} across the 97 communities with at least 40 sales in
2025, the typical miss ranges from 3.9\% (Cedarglen, condos) to 13.5\% (Lindencliffe, 73\% houses);
the middle half is 6.1 to 8.0\%. It is worse where more sales are houses (correlation 0.39) and where prices
are higher (0.43). \emph{Time:} estimates lean about 2\% high in July and August, when the summer dip runs
deeper than projected, and 5\% high in 2022, when prices fell after the spring peak.
\emph{Honest uncertainty:} an 80\% range for each estimate, from the 10th to 90th percentile of the
2023--24 misses for its property type, covered 79.5\% of the 2025 homes, with a typical width of 22.6\% of
the price (about $\pm$11\%). The bootstrap intervals treat homes as independent; homes sold in the same
month share the same market surprise, so the real uncertainty is somewhat wider.

\textbf{What did not help} (all judged on the 2023 and 2024 summers): separate indexes for houses and
condos (the projection noise outweighed the bias it fixed); following community price trends (changes in a
community's premium have autocorrelation $-0.41$, so they reverse rather than persist); blending in a random
forest (9.2\% on its own); separate tree models for houses and condos (within noise); a fixed summer bias correction (the best shift varied from 0 to 1.75\% by year).
\textbf{Limits:} do not trust it for expensive or unusual houses, for thin groups such as co-ops (50 test
homes in 2025), or in a summer when the market turns sharply. Data on age, renovation, floor level and condo
fees would help most.
